\section{Exact Finite CAR Products: Density Waves and Pairing}
\label{sec:ch20-projected-products}

A single fermion field changes particle number by one.  A charge-density-wave
(CDW) operator, by contrast, removes a left-moving fermion and creates a
right-moving fermion; a superconducting (SC) pair-annihilation operator removes
one fermion of each chirality.  Both are products of two fields.  This simple
observation raises a question which does not arise for a single field:
\emph{what happens between the first operation and the second when our
calculation is confined to an energy-budget subspace?}

The full finite fermionic Fock space has an unambiguous answer: multiply the
operators in their stated order.  Budget restrictions introduce a different
operation, however.  The product of \emph{individually projected} fields
passes through a selected intermediate subspace, whereas a projection of the
\emph{full product} allows intermediate states to leave that subspace and
return before the final projection.  Those discarded paths form an explicit
remainder.  Keeping the remainder is the purpose of this section.  The
construction also provides a convenient setting in which to check chirality,
charge labels, Fourier phases, and fermionic signs before introducing the
formal vertex operators of Section~\ref{sec:ch20-formal-vertices}.

Throughout this section, all identities concerning finite fermion operators
are elementary statements on a finite-dimensional Hilbert space; their
application to the chapter's proposed bosonization dictionary remains a
separate mathematical verification task.  In particular, no new Chapter~20
Lean theorem is being claimed here.

\subsection{The ambient finite-CAR system}
\label{subsec:ch20-ambient}

Fix an integer \(h\geq 1\) and put \(L=2h\).  The spatial lattice is
\(\mathbb Z/L\mathbb Z\), represented by integers \(x\) modulo \(L\); the
centered momentum band is
\begin{equation}
  \mathcal B_L=\{-h+1,-h+2,\ldots,h\},\qquad
  \zeta=\exp(2\pi i/L).
  \label{eq:ch20-band}
\end{equation}
There are two distinct species, called right and left and denoted
\(R\) and \(L\) in subscripts.  The letter \(L\) without a subscript
continues to mean lattice length.  We order the single-particle orbitals
\emph{first by species} \(R<L\), then by increasing momentum within each
species.  This order fixes the occupation-basis signs and is not a harmless
convention if intermediate formulas are compared term by term.

Let
\(\mathcal F=\bigwedge(\mathbb C^{\{R,L\}\times\mathcal B_L})\)
be the fermionic Fock Hilbert space.  Its mode annihilators and creators
\(a_{\nu k}\) and \(a_{\nu k}^{\dagger}\) obey the canonical
anticommutation relations (CAR)
\begin{align}
 \{a_{\nu k},a_{\eta q}\}&=0,
 &\{a_{\nu k}^{\dagger},a_{\eta q}^{\dagger}\}&=0,
 &\{a_{\nu k},a_{\eta q}^{\dagger}\}&=
             \delta_{\nu\eta}\delta_{kq}I.
 \label{eq:ch20-car}
\end{align}
Here \(\{A,B\}=AB+BA\), and the rightmost factor acts first.  The
periodic-twist fields used in the calculations below are
\begin{equation}
 c_{\nu,x}=\frac1{\sqrt L}\sum_{k\in\mathcal B_L}
       \zeta^{kx}a_{\nu k},\qquad
 c_{\nu,x}^{\dagger}=\frac1{\sqrt L}\sum_{k\in\mathcal B_L}
       \zeta^{-kx}a_{\nu k}^{\dagger}.
 \label{eq:ch20-fields}
\end{equation}
Thus the annihilation field carries a \emph{positive} Fourier exponent,
whereas the creator carries its complex conjugate.  Both species use this
same computational band; the left-moving physical branch is encoded by the
branch identification rather than by reversing the exponent in
\eqref{eq:ch20-fields}.  In particular, a continuum oscillation
\(e^{\pm2ik_Fx}\) has \emph{not} been inserted by hand.  Formula
\eqref{eq:ch20-fields} is the periodic-twist specialization of the
finite fields from Chapter~14; a general boundary twist contributes
additional factors and is not covered by the phase formulas derived below.

For later comparison with positive-frequency currents, we reserve
\(M\in\mathbb N\) for an optional mode cutoff with labels
\(m=1,\ldots,M\) (an empty set when \(M=0\)).  The finite-CAR fields
\eqref{eq:ch20-fields} sum over the \emph{whole} band \(\mathcal B_L\);
there is no cutoff \(M\) in their definition or in the exact product
identities of this section.  Restrictions on \(M\) needed for a particular
non-wrapping current representation must be imposed in that representation.

\begin{definition}[Finite CDW and pair-annihilation words]
\label{def:ch20-order-words}
At a lattice site \(x\), define bounded endomorphisms of \(\mathcal F\) by
\begin{equation}
  O_{\mathrm{CDW}}(x)=c_{R,x}^{\dagger}c_{L,x},\qquad
  O_{\mathrm{SC}}(x)=c_{R,x}c_{L,x}.
  \label{eq:ch20-order-words}
\end{equation}
The order of the factors in each definition is part of the definition.
\(O_{\mathrm{SC}}\) annihilates a pair; its adjoint
\(O_{\mathrm{SC}}^{\dagger}=c_{L,x}^{\dagger}c_{R,x}^{\dagger}\)
creates a pair in the reversed order dictated by the adjoint.
\end{definition}

Both operators are \emph{even CAR words}: a fermionic parity transformation
changes the sign of each field but not of their product.  Evenness does not
mean charge neutrality.  The CDW operator is neutral only with respect to
the \emph{total} charge, whereas the SC operator lowers it by two.

To state this exactly, let
\begin{equation}
  \widehat n_{\nu}=\sum_{k\in\mathcal B_L}
     a_{\nu k}^{\dagger}a_{\nu k},\qquad
  Q_{\nu}=\widehat n_{\nu}-hI,
 \label{eq:ch20-charges}
\end{equation}
so that the reference half-filled sea carries zero relative charge
\(Q_R=Q_L=0\).  For a charge label
\(N=(N_R,N_L)\in\mathbb Z^2\), let
\(e_R=(1,0)\) and \(e_L=(0,1)\).

\begin{proposition}[Exact charge shifts]
\label{prop:ch20-charge-shifts}
For either \(\nu\in\{R,L\}\), the following operator identities hold on
\emph{all} of \(\mathcal F\):
\begin{align}
 [Q_{\nu},O_{\mathrm{CDW}}(x)]
   &=\bigl(\delta_{\nu R}-\delta_{\nu L}\bigr)
                       O_{\mathrm{CDW}}(x),
 \label{eq:ch20-cdw-charge}\\
 [Q_{\nu},O_{\mathrm{SC}}(x)]
   &=-\bigl(\delta_{\nu R}+\delta_{\nu L}\bigr)
                       O_{\mathrm{SC}}(x).
 \label{eq:ch20-sc-charge}
\end{align}
Consequently, with \(Q_{\mathrm{tot}}=Q_R+Q_L\),
\([Q_{\mathrm{tot}},O_{\mathrm{CDW}}]=0\) and
\([Q_{\mathrm{tot}},O_{\mathrm{SC}}]=-2O_{\mathrm{SC}}\).
\end{proposition}
\begin{proof}
For a single orbital \(j=(\eta,q)\), the CAR imply
\([a_j^{\dagger}a_j,a_i]=-\delta_{ji}a_i\) and
\([a_j^{\dagger}a_j,a_i^{\dagger}]=\delta_{ji}a_i^{\dagger}\).
For example,
\(a_j^{\dagger}a_ja_i-a_i a_j^{\dagger}a_j
 =-a_j^{\dagger}a_i a_j-(\delta_{ij}-a_j^{\dagger}a_i)a_j
 =-\delta_{ij}a_j\), which equals \(-\delta_{ij}a_i\).
Summing over \(q\) and using the linear Fourier transforms gives
\begin{equation}
 [Q_\nu,c_{\eta,x}]=-\delta_{\nu\eta}c_{\eta,x},
 \qquad
 [Q_\nu,c_{\eta,x}^{\dagger}]
       =\delta_{\nu\eta}c_{\eta,x}^{\dagger}.
 \label{eq:ch20-single-charge}
\end{equation}
The scalar \(-hI\) commutes with every operator.  Finally apply the
commutator product rule \([Q,AB]=[Q,A]B+A[Q,B]\) to each ordered
product in \eqref{eq:ch20-order-words}; summing the two charge identities
proves the final assertions.
\end{proof}

In particular, the source, intermediate and target charges for our two
products are forced by the operators, not chosen independently:
\begin{equation}
\begin{aligned}
 N_s=N\quad&\xrightarrow{\ c_{L,x}\ }
       N_m=N-e_L\\[-2pt]
       &\xrightarrow{\ c_{R,x}^{\dagger}\ }
       N_t^{\mathrm{CDW}}=N+e_R-e_L,\\[3pt]
 N_s=N\quad&\xrightarrow{\ c_{L,x}\ }
       N_m=N-e_L\\[-2pt]
       &\xrightarrow{\ c_{R,x}\ }
       N_t^{\mathrm{SC}}=N-e_R-e_L.
\end{aligned}
\label{eq:ch20-charge-arrows}
\end{equation}
The arrows are drawn in \emph{time-of-application order}, from right to
left in the written CAR word.  Because \(c_R\) and \(c_L\)
anticommute, reversing their positions introduces a minus sign even
though the charge shift is unchanged.

\subsection{Three budget carriers, not one}
\label{subsec:ch20-budget}

For each admissible charge \(N_\nu\in[-h,h]\cap\mathbb Z\), the
\emph{sector ground occupation} fills momenta
\(\{-h+1,\ldots,N_\nu\}\) for species \(\nu\); if
\(N_\nu=-h\) the set is empty.  Write \(|\Omega_N\rangle\) for the
normalized occupation ket obtained by creating all these particles in
the fixed \(R<L\), increasing-momentum order.  Thus
\begin{equation}
  Q_\nu|\Omega_N\rangle=N_\nu|\Omega_N\rangle.
 \label{eq:ch20-ground-charge}
\end{equation}
An occupation configuration \(S=(S_R,S_L)\) in this sector has a
nonnegative integer excitation grade
\begin{equation}
 E_N(S)=\sum_{\nu=R,L}\left(
      \sum_{k\in S_\nu}k
      -\sum_{k=-h+1}^{N_\nu}k\right).
 \label{eq:ch20-excitation}
\end{equation}
Indeed, among subsets of \(\mathcal B_L\) of cardinality \(h+N_\nu\),
the lowest \(h+N_\nu\) momenta minimize the sum.  Equivalently,
\(E_N(S)\) counts the net upward displacement of particles relative to
the sector ground.  The same grading underlies the budget spaces of
Chapter~\ref{ch:vacuum_budget}; in the present finite-dimensional carrier
it can be read directly from occupation kets.

For \(K\in\mathbb N\), define the cutoff carrier
\begin{equation}
 V_{N,K}=\operatorname{span}\{\,|S\rangle:
       Q_\nu|S\rangle=N_\nu|S\rangle\ (\nu=R,L),
       \ E_N(S)\leq K\,\}\ \subseteq\mathcal F.
 \label{eq:ch20-budget-carrier}
\end{equation}
For an inadmissible charge we take \(V_{N,K}=\{0\}\).  Each such space
has the inner product inherited from \(\mathcal F\).  Denote its
isometric inclusion by \(i_{N,K}:V_{N,K}\hookrightarrow\mathcal F\)
and its orthogonal projection by
\(p_{N,K}:\mathcal F\to V_{N,K}\).  The associated ambient
projection is
\begin{equation}
 P_{N,K}=i_{N,K}p_{N,K},\quad
 P_{N,K}^{2}=P_{N,K}=P_{N,K}^{\dagger},\quad
 p_{N,K}i_{N,K}=I_{V_{N,K}}.
 \label{eq:ch20-typed-projections}
\end{equation}
The notation distinguishes an endomorphism of the ambient Fock space
from a map between two \emph{different} budget carriers.  These are not
interchangeable objects, even if they happen to have the same dimension.

Choose three cutoff integers \(K_s,K_m,K_t\geq0\).  In the formulas
below, the abbreviations
\begin{equation}
\begin{gathered}
 V_s=V_{N,K_s},\qquad V_m=V_{N-e_L,K_m},\\
 V_t^{J}=V_{N_t^{J},K_t},\qquad
 i_s=i_{N,K_s},\quad i_m=i_{N-e_L,K_m},\quad
 p_m=p_{N-e_L,K_m},\quad P_m=i_mp_m
\end{gathered}
\label{eq:ch20-three-carriers}
\end{equation}
are \emph{typed} by the charge labels in
\eqref{eq:ch20-charge-arrows}; the target projection is written
\(p_t^J=p_{N_t^J,K_t}\).  No condition such as
\(K_s=K_m=K_t\) is imposed.  The intermediate cutoff is a separate
choice and must not be silently identified with the target cutoff.

\begin{definition}[Restricted one-field arrows]
\label{def:ch20-field-arrows}
Set
\begin{equation}
 B_L=p_m c_{L,x}i_s:V_s\longrightarrow V_m,
 \label{eq:ch20-BL}
\end{equation}
and define the second arrows, with the indicated \emph{different targets},
by
\begin{align}
 B_R^{\mathrm{CDW},\dagger}
   &=p_t^{\mathrm{CDW}}c_{R,x}^{\dagger}i_m:
         V_m\longrightarrow V_t^{\mathrm{CDW}},
 \label{eq:ch20-BR-CDW}\\
 B_R^{\mathrm{SC}}
   &=p_t^{\mathrm{SC}}c_{R,x}i_m:
         V_m\longrightarrow V_t^{\mathrm{SC}}.
 \label{eq:ch20-BR-SC}
\end{align}
The dagger in \(B_R^{\mathrm{CDW},\dagger}\) refers to the underlying
ambient creator.  It does \emph{not} purport to take the adjoint of
\(B_L\), nor of a field with unrelated source or target data.
\end{definition}

There is a genuine Hilbert-adjoint statement behind the creator notation.
Define the oppositely directed map
\(D_R=p_m c_{R,x}i_{N_t^{\mathrm{CDW}},K_t}:
V_t^{\mathrm{CDW}}\to V_m\).  Then
\begin{equation}
 D_R^{\dagger}
   =p_t^{\mathrm{CDW}}c_{R,x}^{\dagger}i_m
   =B_R^{\mathrm{CDW},\dagger}.
 \label{eq:ch20-reversed-adjoint}
\end{equation}
To verify this rather than infer it from notation, take
\(u\in V_m\) and \(v\in V_t^{\mathrm{CDW}}\):
\begin{align*}
 \langle D_Rv,u\rangle_{V_m}
 &=\langle c_{R,x}i_tv,i_mu\rangle_{\mathcal F}\\
 &=\langle i_tv,c_{R,x}^{\dagger}i_mu\rangle_{\mathcal F}
  =\langle v,p_t^{\mathrm{CDW}}c_{R,x}^{\dagger}i_mu\rangle_{V_t^{\mathrm{CDW}}}.
\end{align*}
Here \(i_t=i_{N_t^{\mathrm{CDW}},K_t}\) in this display.
Thus the adjoint reverses \emph{both} the charge data and the cutoffs.
This is the precise finite Hilbert-space meaning of Chapter~14's
budget-adjoint construction.

\begin{theorem}[Two-field projection insertion and retained residual]
\label{thm:ch20-projection-residual}
For every admissible or inadmissible \(N\), every choice of
\(K_s,K_m,K_t\geq0\), and every \(x\in\mathbb Z\), the following
identities of linear maps hold with the indicated source and target:
\begin{align}
 p_t^{\mathrm{CDW}}O_{\mathrm{CDW}}(x)i_s
 &=B_R^{\mathrm{CDW},\dagger}\circ B_L
   +\mathcal R_{\mathrm{CDW}}^{smt}(x),
 &V_s&\longrightarrow V_t^{\mathrm{CDW}},
 \label{eq:ch20-cdw-residual}\\
 p_t^{\mathrm{SC}}O_{\mathrm{SC}}(x)i_s
 &=B_R^{\mathrm{SC}}\circ B_L
   +\mathcal R_{\mathrm{SC}}^{smt}(x),
 &V_s&\longrightarrow V_t^{\mathrm{SC}},
 \label{eq:ch20-sc-residual}
\end{align}
where the \emph{uncompressed intermediate excursions} are
\begin{align}
 \mathcal R_{\mathrm{CDW}}^{smt}(x)
   &=p_t^{\mathrm{CDW}}c_{R,x}^{\dagger}(I_{\mathcal F}-P_m)
                                   c_{L,x}i_s,
 \label{eq:ch20-cdw-R}\\
 \mathcal R_{\mathrm{SC}}^{smt}(x)
   &=p_t^{\mathrm{SC}}c_{R,x}(I_{\mathcal F}-P_m)
                                   c_{L,x}i_s.
 \label{eq:ch20-sc-R}
\end{align}
In particular, the factorization into the two projected single fields
is exact \emph{if and only if} the corresponding displayed residual map
vanishes.  No such vanishing is assumed in the identities.
\end{theorem}

\begin{proof}
The ambient projection satisfies the tautological decomposition
\(I_{\mathcal F}=P_m+(I_{\mathcal F}-P_m)\).
Insert it between the two factors of the CDW word:
\begin{align*}
 p_t^{\mathrm{CDW}}c_{R,x}^{\dagger}c_{L,x}i_s
 & =p_t^{\mathrm{CDW}}c_{R,x}^{\dagger}P_m c_{L,x}i_s
   +p_t^{\mathrm{CDW}}c_{R,x}^{\dagger}
                        (I_{\mathcal F}-P_m)c_{L,x}i_s\\
 & =\bigl(p_t^{\mathrm{CDW}}c_{R,x}^{\dagger}i_m\bigr)
    \circ\bigl(p_m c_{L,x}i_s\bigr)
   +\mathcal R_{\mathrm{CDW}}^{smt}(x),
\end{align*}
where the second equality uses \(P_m=i_mp_m\).  This is
\eqref{eq:ch20-cdw-residual}.  Replace the first creator by
\(c_{R,x}\) to obtain \eqref{eq:ch20-sc-residual}.  All expressions
are maps between the spaces explicitly shown in
\eqref{eq:ch20-three-carriers}; the proof makes no assumption about
how many intermediate excitations a field can generate.
Finally, subtract the factored term from either equality: it vanishes
exactly when its residual does.
\end{proof}

An immediate \emph{sufficient}, though not necessary, vanishing
criterion is
\begin{equation}
 (I_{\mathcal F}-P_m)c_{L,x}i_s=0,
 \quad\text{equivalently}\quad
 c_{L,x}(V_s)\subseteq i_m(V_m).
 \label{eq:ch20-safe-intermediate}
\end{equation}
This is the strongest, easiest-to-check form of an intermediate
\emph{range condition}: every state produced by the first field already
lies in the proposed intermediate budget.  A weaker possibility is that
the second field and final projection kill all excluded states:
\(p_t^J c_{R,x}^{\sharp}(I-P_m)c_{L,x}i_s=0\), where
\(c_{R,x}^{\sharp}=c_{R,x}^{\dagger}\) for CDW and
\(c_{R,x}^{\sharp}=c_{R,x}\) for SC.  Such a cancellation is about the
\emph{complete} product, and cannot be inferred by checking each factor
in isolation.

\begin{remark}[Why an intermediate projection changes multiplication]
\label{rem:ch20-not-multiplicative}
Compression \(A\mapsto p A i\) is linear, but generally not
multiplicative.  An intermediate state can lie outside \(V_m\) and still
be mapped into \(V_t^J\) by the second field.  Inserting \(P_m\)
deletes that route.  Even a correct identity for each \emph{one-field}
projected map therefore does not license multiplying those identities
and discarding \(\mathcal R_J^{smt}\).  In particular, Chapter~14's
finite Taylor candidates, if separately shown to equal the appropriate
individual projected fields under its conditional cyclic criterion,
still require the present residual and the actual transition cutoffs when
composed.  No factorization, commutation, or Baker--Campbell--Hausdorff
identity for \emph{compressed exponentials} follows from
\eqref{eq:ch20-cdw-residual}--\eqref{eq:ch20-sc-residual}.
\end{remark}

\subsection{Charge phases and the ordered Klein signs}
\label{subsec:ch20-klein}

The operator identities above require no vacuum or special reference
state.  It is nevertheless useful to compute the matrix elements between
\emph{sector ground kets}, both as a stringent sign check and as the
zero-excitation portion of a possible order-parameter dictionary.  We now
assume the periodic twist in \eqref{eq:ch20-fields} and the fixed
occupation ordering \(R<L\).

The normalized occupation ket at charge \(N\) contains exactly
\(h+N_\nu\) particles of species \(\nu\).  In the occupation basis,
annihilation of an occupied orbital contributes \((-1)^n\), with \(n\)
the number of occupied orbitals preceding it; creation of an unoccupied
orbital obeys the same sign rule.  Define the signed Klein coefficient for
removal of the \emph{highest occupied} orbital of a species by
\begin{equation}
 \begin{aligned}
 P_R(N)&=(-1)^{h+N_R-1},\\
 P_L(N)&=(-1)^{(h+N_R)+(h+N_L-1)}.
 \end{aligned}
 \label{eq:ch20-klein-phases}
\end{equation}
Equivalently,
\begin{equation}
 P_\nu(N)=(-1)^{\displaystyle
     \sum_{\eta<\nu}(h+N_\eta)+h+N_\nu-1},
 \label{eq:ch20-klein-general}
\end{equation}
where \(\eta<\nu\) refers to the fixed species ordering and
\((-1)^n\) is defined for every integer \(n\) by its parity.
The formula is the charge-sector sign convention of Chapter~13,
not an independently adjustable choice of relative Klein phases.

For removal of the last occupied momentum \(k=N_\nu\), the annihilation
field supplies the Fourier coefficient \(L^{-1/2}\zeta^{N_\nu x}\).
For addition of the first unoccupied momentum \(k=N_R+1\), the creator
supplies \(L^{-1/2}\zeta^{-(N_R+1)x}\).  The creation sign equals
\(P_R(N+e_R)\): it counts \(h+N_R\) occupied right orbitals
\emph{before} the new particle is added.  This shifted argument matters.
Likewise, after the left annihilation in an SC product the right removal
sign is \(P_R(N-e_L)\); no factor may be moved past the other without
recomputing its source label.

\begin{proposition}[Exact sector-ground matrix elements]
\label{prop:ch20-ground-coeff}
Let \(-h\leq N_R,N_L\leq h\), and let
\(\langle\Omega_{N'}|A|\Omega_N\rangle\) denote the finite-Fock
Hilbert-space matrix element.  If \(N_L>-h\) and \(N_R<h\), then
\begin{align}
 &\langle\Omega_{N+e_R-e_L}|
                O_{\mathrm{CDW}}(x)|\Omega_N\rangle\\
 &\hspace{1cm}=\frac1L
     \zeta^{(N_L-N_R-1)x}
     P_L(N)P_R(N-e_L+e_R).
 \label{eq:ch20-cdw-ground}
\end{align}
If \(N_L>-h\) and \(N_R>-h\), then
\begin{align}
 &\langle\Omega_{N-e_R-e_L}|
                O_{\mathrm{SC}}(x)|\Omega_N\rangle\\
 &\hspace{1cm}=\frac1L
     \zeta^{(N_R+N_L)x}
     P_L(N)P_R(N-e_L).
 \label{eq:ch20-sc-ground}
\end{align}
If a required removal or creation is Pauli-forbidden at the corresponding
empty or full boundary, the respective ground-to-ground matrix element is
zero.  The nonzero coefficients above have modulus \(1/L\).
\end{proposition}
\begin{proof}
Expand the product using \eqref{eq:ch20-fields}:
\begin{align*}
 O_{\mathrm{CDW}}(x)
  &=\frac1L\sum_{p,q\in\mathcal B_L}
      \zeta^{(-p+q)x}a_{Rp}^{\dagger}a_{Lq},\\
 O_{\mathrm{SC}}(x)
  &=\frac1L\sum_{p,q\in\mathcal B_L}
      \zeta^{(p+q)x}a_{Rp}a_{Lq}.
\end{align*}
In the CDW case the target ket contains precisely the occupied left
orbitals \(-h+1,\ldots,N_L-1\); hence the \emph{only} term that
can connect the two specified ground configurations has \(q=N_L\).
Similarly, the target right orbitals are \(-h+1,\ldots,N_R+1\),
which forces \(p=N_R+1\).  The left field acts first, giving
\(P_L(N)\).  The right creator then adds the first empty right
orbital to the intermediate state \(N-e_L\), with sign
\((-1)^{h+N_R}=P_R(N-e_L+e_R)\).  The two Fourier characters
multiply to \(\zeta^{(N_L-N_R-1)x}\).  All other terms are either
zero by Pauli exclusion or produce an occupation ket orthogonal to
the stated target ground.

For SC the target has one fewer particle of each species.  Its ground
occupation forces \(q=N_L\) and \(p=N_R\).  Acting first with the
left field yields \(P_L(N)\); the subsequent right removal sees the
charge \(N-e_L\), and gives \(P_R(N-e_L)\).  Its Fourier factor is
\(\zeta^{(N_L+N_R)x}\).  The asserted boundary cases follow because
a missing highest occupied momentum cannot be annihilated, and a full
right band offers no unoccupied orbital for CDW creation.  Because
\(|\zeta|=|P_\nu|=1\), the modulus of each allowed matrix element
is \(1/L\).
\end{proof}

It is helpful to separate the phase coming from the \emph{source
occupation} from the signed shift.  For these ground matrix elements
write
\begin{align}
 \chi_{\mathrm{CDW}}(N,x)&=\zeta^{(N_L-N_R-1)x},
 &\kappa_{\mathrm{CDW}}(N)
   &=P_L(N)P_R(N-e_L+e_R),
 \label{eq:ch20-cdw-chi-kappa}\\
 \chi_{\mathrm{SC}}(N,x)&=\zeta^{(N_R+N_L)x},
 &\kappa_{\mathrm{SC}}(N)
   &=P_L(N)P_R(N-e_L).
 \label{eq:ch20-sc-chi-kappa}
\end{align}
The scalar \(\chi_J(N,x)\) acts on the \emph{source} charge label
before either shift is performed.  The left removal acts first in
both products.  Its Klein sign is evaluated at \(N\); the next
Klein sign uses the appropriate shifted charge.  These functions are
\emph{matrix-element coefficients}; they are not an assertion that
\(O_J\) equals a charge shift multiplied by a scalar on arbitrary
excited states.  Excited states require the full CAR word and, when
budget-restricted, its residual.

\subsection{A four-site calculation: the residual is genuinely present}
\label{subsec:ch20-four-site}

\begin{example}[Both order parameters and a missed intermediate path]
\label{ex:ch20-four-site}
Take \(h=2\), hence \(L=4\),
\(\mathcal B_4=\{-1,0,1,2\}\),
\(\zeta=i\), and \(x=1\).  For compactness write
\(|R\{\ldots\};L\{\ldots\}\rangle\) for the normalized occupation
ket with the indicated occupied momenta, in the \(R<L\) ordering.

\emph{First: the CDW ground coefficient.}
The source at \(N=(0,0)\) is
\[
 |\Omega_{(0,0)}\rangle
  =|R\{-1,0\};L\{-1,0\}\rangle.
\]
After a left removal and right creation the target sector is
\(N_t=(1,-1)\), with ground
\[
 |\Omega_{(1,-1)}\rangle
  =|R\{-1,0,1\};L\{-1\}\rangle.
\]
Only \(a_{R,1}^{\dagger}a_{L,0}\) connects these two ground kets.
The Fourier product is \(\zeta^{-1}= -i\), the left removal sign is
\(P_L(0,0)=(-1)^3=-1\), and the right creation sign is
\(P_R(1,-1)=(-1)^2=+1\).  Therefore
\begin{equation}
 \langle\Omega_{(1,-1)}|
    O_{\mathrm{CDW}}(1)|\Omega_{(0,0)}\rangle
 =\frac14(-i)(-1)(+1)=\frac{i}{4}.
 \label{eq:ch20-four-cdw}
\end{equation}

\emph{Second: the SC ground coefficient.}
Choose the source \(N=(1,0)\):
\[
 |\Omega_{(1,0)}\rangle
   =|R\{-1,0,1\};L\{-1,0\}\rangle.
\]
The target charge is \((0,-1)\), with ground
\( |\Omega_{(0,-1)}\rangle
   =|R\{-1,0\};L\{-1\}\rangle\).
Only \(a_{R,1}a_{L,0}\) connects these states.  Its Fourier
character is \(\zeta^{1+0}=i\), and the two successive signs are
\(P_L(1,0)=+1\) and \(P_R(1,-1)=+1\).  It follows that
\begin{equation}
 \langle\Omega_{(0,-1)}|
    O_{\mathrm{SC}}(1)|\Omega_{(1,0)}\rangle
 =\frac{i}{4}.
 \label{eq:ch20-four-sc}
\end{equation}
These two coefficients are exact finite-Fock numbers, not
large-system approximations or formal vertex-state evaluations.

\emph{Third: a discarded intermediate excitation survives the
second field.}
Return to the CDW source \(N=(0,0)\) and select
\(K_s=K_m=0\), \(K_t=1\).  The intermediate charge is
\(N_m=(0,-1)\) and its zero-budget state has occupied orbitals
\(R\{-1,0\},L\{-1\}\).  Consider instead the term in which
\(c_{L,1}\) removes the \emph{left} particle at \(q=-1\):
\[
 |R\{-1,0\};L\{-1,0\}\rangle
 \xrightarrow{\ a_{L,-1}\ }
 |R\{-1,0\};L\{0\}\rangle.
\]
The latter configuration has the correct intermediate \emph{charge}
\((0,-1)\), but its left particle occupies \(0\) rather than \(-1\).
By \eqref{eq:ch20-excitation} it has excitation \(E_{N_m}=1\), and
is therefore killed by \(P_m=P_{(0,-1),0}\).  The right creator can
nevertheless add \(p=1\):
\[
 |R\{-1,0\};L\{0\}\rangle
 \xrightarrow{\ a_{R,1}^{\dagger}\ }
 |R\{-1,0,1\};L\{0\}\rangle.
\]
This final state belongs to \(V_{(1,-1),1}\), having left excitation
one and right excitation zero.  With the stated ordering, both mode
operations have occupation sign \(+1\); their Fourier factors are
\(\zeta^{-1}=-i\) for the left annihilator and
\(\zeta^{-1}=-i\) for the right creator.  No other excluded
intermediate path contributes to this final occupation ket.  Hence
\begin{equation}
 \begin{aligned}
 &\mathcal R_{\mathrm{CDW}}^{smt}(1)
                |\Omega_{(0,0)}\rangle\\
 &\qquad=-\frac14
      |R\{-1,0,1\};L\{0\}\rangle\ne0
       \quad (K_s=K_m=0,\ K_t=1).
 \end{aligned}
 \label{eq:ch20-four-residual}
\end{equation}
In particular, equating the projected product with
\(B_R^{\mathrm{CDW},\dagger}B_L\) in this example would
be \emph{false}, despite the fact that both constituent fields and
all three projections are correctly defined.

If we change only the \emph{target} cutoff to \(K_t=0\), the state
on the right of \eqref{eq:ch20-four-residual} is also discarded.
There is then no surviving excluded path from this source ground to
the target ground, and
\begin{equation}
 \mathcal R_{\mathrm{CDW}}^{smt}(1)
      |\Omega_{(0,0)}\rangle=0
 \qquad(K_s=K_m=K_t=0).
 \label{eq:ch20-four-residual-zero}
\end{equation}
The restricted, ground-to-ground equality here is a property of the
\emph{final projection} and this source ket; it does not turn
\(\mathcal R_{\mathrm{CDW}}^{smt}\) into the zero operator for
arbitrary larger budgets.
\end{example}

The example illustrates two conceptually different uses of a cutoff.
The intermediate budget \(K_m=0\) declares an excited left-hole state
inadmissible \emph{before} the second field acts.  The target budget
\(K_t=1\) still accepts its eventual image.  This mismatch creates the
nonzero remainder; it is not an error in the CAR algebra or the Fourier
normalization.  Conversely, if one is interested only in a particular
ground-to-ground coefficient, orthogonality of the untouched left
occupation can eliminate the excluded paths \emph{after} the full
product.  Neither conclusion justifies a blanket rule that energy-budget
projection commutes with multiplication.

\subsection{What these exact results do and do not provide}
\label{subsec:ch20-scope}

The finite CAR construction now supplies three distinct pieces of data:
a globally defined even word \(O_J(x)\); a collection of typed maps
\(V_s\to V_t^J\) with explicit intermediate remainder
\(\mathcal R_J^{smt}\); and a charge-ground scalar coefficient with
specified Fourier and Klein ordering.  These statements require no
choice of finite interacting vacuum and no canonical-commutation
relations for compressed density modes.

For a normalized \(\psi\in\mathcal F\), an observable expectation such
as \(\langle\psi,O_J(x)^{\dagger}O_J(y)\psi\rangle\) is an ordinary
finite Hilbert-space number.  A density matrix \(\rho\geq0\),
\(\operatorname{Tr}\rho=1\), similarly defines
\(\operatorname{Tr}(\rho O_J(x)^{\dagger}O_J(y))\).  These should
not be identified without a comparison theorem with the abstract
positive functional on the \emph{uncompressed} oscillator algebra
discussed in Chapter~\ref{ch:correlations} and the construction
underlying Section~\ref{sec:ch20-formal-vertices}.  In particular,
a formal vertex correlator evaluated coefficientwise in a power-series
parameter is not automatically the finite-CAR expectation just defined.

A proposed Mattis--Mandelstam description may eventually replace a
single projected field by a finite polynomial or oscillator expression,
but Chapter~14's conditional criterion must first be checked for the
\emph{specified source and target cutoffs}.  The exact product theorem
then retains \(\mathcal R_J^{smt}\) unless an explicit range or
final-projection annihilation argument removes it.  This distinction is
especially important for later pairing, duality and spinful/Umklapp
constructions: a scalar agreement on sector grounds does not certify an
operator identity on all excited budget states.

Finally, the sign and phase conventions fixed here are provisional
cross-chapter conventions: \(R<L\), common positive-exponent finite
annihilation Fourier transform, left field acting first, source charge
phases, and separate cutoffs for each arrow.  Their compatibility with
future chapter integrations must be reviewed against the actual committed
interfaces.  Neither a thermodynamic limit nor continuum scaling,
spectral equivalence or a finite interacting-state identification has
been established in this section.
